
\documentclass{article}
\usepackage{colortbl}
\usepackage{makecell}
\usepackage{multirow}
\usepackage{supertabular}

\begin{document}

\newcounter{utterance}

\centering \large Interaction Transcript for game `dond', experiment `semi', episode 3 with gpt{-}5.4{-}azure.
\vspace{24pt}

{ \footnotesize  \setcounter{utterance}{1}
\setlength{\tabcolsep}{0pt}
\begin{supertabular}{c@{$\;$}|p{.15\linewidth}@{}p{.15\linewidth}p{.15\linewidth}p{.15\linewidth}p{.15\linewidth}p{.15\linewidth}}
    \# & $\;$A & \multicolumn{4}{c}{Game Master} & $\;\:$B\\
    \hline

    \theutterance \stepcounter{utterance}  
    & & \multicolumn{4}{p{0.6\linewidth}}{
        \cellcolor[rgb]{0.9,0.9,0.9}{
            \makecell[{{p{\linewidth}}}]{
                \texttt{\tiny{[P1$\langle$GM]}}
                \texttt{Grasz w grę negocjacji kooperacyjnej, w której musisz uzgodnić z innym graczem sposób podziału zestawu przedmiotów między sobą.} \\
\\ 
\texttt{Zasady:} \\
\texttt{(a) Ty i drugi gracz otrzymujecie zestaw przedmiotów. Każdemu z Was przydzielana jest również tajna funkcja wartości, reprezentująca, ile warte są dla Ciebie poszczególne typy przedmiotów.} \\
\texttt{(b) Wymieniasz wiadomości z drugim graczem, aby uzgodnić, kto otrzyma jakie przedmioty. Każdy może wysłać maksymalnie 5 wiadomości lub skończyć grę wcześniej, składając tajną propozycję w dowolnym momencie.} \\
\texttt{(c) Poproszony o złożenie tajnej propozycji wskazującej przedmioty, które chcesz, sformatowanej w nawiasach kwadratowych w następujący sposób: "{[}Propozycja: <liczba> <nazwa przedmiotu>, <liczba> <nazwa przedmiotu>, <...>{]}"} \\
\texttt{(d) Jeśli Twoje propozycje są się się nawzajem uzupełniają, tzn. jest wystarczająco przedmiotów, aby spełnić obie propozycje, każdy gracz otrzymuje wynik na podstawie sumy wartości przedmiotów, które otrzymał. W przeciwnym razie obaj otrzymujecie zero punktów.} \\
\texttt{(e) Twój cel to maksymalizacja sumy Twojego wyniku i wyniku drugiego gracza.} \\
\\ 
\texttt{Zacznijmy.} \\
\\ 
\texttt{Zestaw dostępnych przedmiotów to:} \\
\\ 
\texttt{1 obcęgi, 2 cegły, 1 pierścionek, 1 zegarek} \\
\\ 
\texttt{Twoja tajna funkcja wartości to:} \\
\\ 
\texttt{obcęgi: 1, cegła: 0, pierścionek: 2, zegarek: 7} \\
\\ 
\texttt{WAŻNE: Twoje wiadomości, chyba że jest to tajna propozycja, są przesyłane bezpośrednio drugiemu graczowi, więc nie dodawaj żadnej odpowiedzi na zasady ani tekstu ogłaszającego Twoją wiadomość. Aby złożyć tajną propozycję, użyj wskazanego formatu. Nie używaj nawiasów kwadratowych podczas komunikowania się z drugim graczem, w przeciwnym razie zostanie to zinterpretowane jako Twoja tajna propozycja.} \\
            }
        }
    }
    & \\ \\

    \theutterance \stepcounter{utterance}  
    & \multicolumn{4}{p{0.6\linewidth}}{
        \cellcolor[rgb]{0.9,0.9,0.9}{
            \makecell[{{p{\linewidth}}}]{
                \texttt{\tiny{[P1$\rangle$GM]}}
                \texttt{Najbardziej zależy mi na zegarku; cegły są mi praktycznie obojętne, pierścionek ma dla mnie umiarkowaną wartość, a obcęgi niewielką. Jakie są Twoje priorytety, żebyśmy mogli ustawić podział możliwie najlepiej dla nas obojga?} \\
            }
        }
    }
    & & \\ \\

    \theutterance \stepcounter{utterance}  
    & & & \multicolumn{2}{p{0.3\linewidth}}{
        \cellcolor[rgb]{0.9,0.9,0.9}{
            \makecell[{{p{\linewidth}}}]{
                \texttt{\tiny{[GM$|$GM]}}
                \texttt{continue} \\
            }
        }
    }
    & & \\ \\

    \theutterance \stepcounter{utterance}  
    & & & \multicolumn{2}{p{0.3\linewidth}}{
        \cellcolor[rgb]{0.9,0.9,0.9}{
            \makecell[{{p{\linewidth}}}]{
                \texttt{\tiny{[GM$|$GM]}}
                \texttt{Najbardziej zależy mi na zegarku; cegły są mi praktycznie obojętne, pierścionek ma dla mnie umiarkowaną wartość, a obcęgi niewielką. Jakie są Twoje priorytety, żebyśmy mogli ustawić podział możliwie najlepiej dla nas obojga?} \\
            }
        }
    }
    & & \\ \\

    \theutterance \stepcounter{utterance}  
    & & \multicolumn{4}{p{0.6\linewidth}}{
        \cellcolor[rgb]{0.9,0.9,0.9}{
            \makecell[{{p{\linewidth}}}]{
                \texttt{\tiny{[GM$\rangle$P2]}}
                \texttt{Grasz w grę negocjacji kooperacyjnej, w której musisz uzgodnić z innym graczem sposób podziału zestawu przedmiotów między sobą.} \\
\\ 
\texttt{Zasady:} \\
\texttt{(a) Ty i drugi gracz otrzymujecie zestaw przedmiotów. Każdemu z Was przydzielana jest również tajna funkcja wartości, reprezentująca, ile warte są dla Ciebie poszczególne typy przedmiotów.} \\
\texttt{(b) Wymieniasz wiadomości z drugim graczem, aby uzgodnić, kto otrzyma jakie przedmioty. Każdy może wysłać maksymalnie 5 wiadomości lub skończyć grę wcześniej, składając tajną propozycję w dowolnym momencie.} \\
\texttt{(c) Poproszony o złożenie tajnej propozycji wskazującej przedmioty, które chcesz, sformatowanej w nawiasach kwadratowych w następujący sposób: "{[}Propozycja: <liczba> <nazwa przedmiotu>, <liczba> <nazwa przedmiotu>, <...>{]}"} \\
\texttt{(d) Jeśli Twoje propozycje są się się nawzajem uzupełniają, tzn. jest wystarczająco przedmiotów, aby spełnić obie propozycje, każdy gracz otrzymuje wynik na podstawie sumy wartości przedmiotów, które otrzymał. W przeciwnym razie obaj otrzymujecie zero punktów.} \\
\texttt{(e) Twój cel to maksymalizacja sumy Twojego wyniku i wyniku drugiego gracza.} \\
\\ 
\texttt{Zacznijmy.} \\
\\ 
\texttt{Zestaw dostępnych przedmiotów to:} \\
\\ 
\texttt{1 obcęgi, 2 cegły, 1 pierścionek, 1 zegarek} \\
\\ 
\texttt{Twoja tajna funkcja wartości to:} \\
\\ 
\texttt{obcęgi: 1, cegła: 1, pierścionek: 4, zegarek: 3} \\
\\ 
\texttt{WAŻNE: Twoje wiadomości, chyba że jest to tajna propozycja, są przesyłane bezpośrednio drugiemu graczowi, więc nie dodawaj żadnej odpowiedzi na zasady ani tekstu ogłaszającego Twoją wiadomość. Aby złożyć tajną propozycję, użyj wskazanego formatu. Nie używaj nawiasów kwadratowych podczas komunikowania się z drugim graczem, w przeciwnym razie zostanie to zinterpretowane jako Twoja tajna propozycja.} \\
\\ 
\texttt{Najbardziej zależy mi na zegarku; cegły są mi praktycznie obojętne, pierścionek ma dla mnie umiarkowaną wartość, a obcęgi niewielką. Jakie są Twoje priorytety, żebyśmy mogli ustawić podział możliwie najlepiej dla nas obojga?} \\
            }
        }
    }
    & \\ \\

    \theutterance \stepcounter{utterance}  
    & & & \multicolumn{4}{p{0.6\linewidth}}{
        \cellcolor[rgb]{0.9,0.9,0.9}{
            \makecell[{{p{\linewidth}}}]{
                \texttt{\tiny{[GM$\langle$P2]}}
                \texttt{U mnie najwyższy priorytet ma pierścionek, zegarek też jest dość cenny, a obcęgi i cegły mają niską wartość. Jeśli dla Ciebie zegarek jest najważniejszy, to sensownie byłoby ustawić to tak, że Ty bierzesz zegarek, ja pierścionek, a cegły i obcęgi możemy dorzucić tam, gdzie bardziej Ci pasują.} \\
            }
        }
    }
     \\ \\

    \theutterance \stepcounter{utterance}  
    & & & \multicolumn{2}{p{0.3\linewidth}}{
        \cellcolor[rgb]{0.9,0.9,0.9}{
            \makecell[{{p{\linewidth}}}]{
                \texttt{\tiny{[GM$|$GM]}}
                \texttt{continue} \\
            }
        }
    }
    & & \\ \\

    \theutterance \stepcounter{utterance}  
    & & & \multicolumn{2}{p{0.3\linewidth}}{
        \cellcolor[rgb]{0.9,0.9,0.9}{
            \makecell[{{p{\linewidth}}}]{
                \texttt{\tiny{[GM$|$GM]}}
                \texttt{U mnie najwyższy priorytet ma pierścionek, zegarek też jest dość cenny, a obcęgi i cegły mają niską wartość. Jeśli dla Ciebie zegarek jest najważniejszy, to sensownie byłoby ustawić to tak, że Ty bierzesz zegarek, ja pierścionek, a cegły i obcęgi możemy dorzucić tam, gdzie bardziej Ci pasują.} \\
            }
        }
    }
    & & \\ \\

    \theutterance \stepcounter{utterance}  
    & & \multicolumn{4}{p{0.6\linewidth}}{
        \cellcolor[rgb]{0.9,0.9,0.9}{
            \makecell[{{p{\linewidth}}}]{
                \texttt{\tiny{[P1$\langle$GM]}}
                \texttt{U mnie najwyższy priorytet ma pierścionek, zegarek też jest dość cenny, a obcęgi i cegły mają niską wartość. Jeśli dla Ciebie zegarek jest najważniejszy, to sensownie byłoby ustawić to tak, że Ty bierzesz zegarek, ja pierścionek, a cegły i obcęgi możemy dorzucić tam, gdzie bardziej Ci pasują.} \\
            }
        }
    }
    & \\ \\

    \theutterance \stepcounter{utterance}  
    & \multicolumn{4}{p{0.6\linewidth}}{
        \cellcolor[rgb]{0.9,0.9,0.9}{
            \makecell[{{p{\linewidth}}}]{
                \texttt{\tiny{[P1$\rangle$GM]}}
                \texttt{Brzmi dobrze. Żeby domknąć to bez ryzyka konfliktu propozycji, proponuję dokładnie tak: ja biorę zegarek i obcęgi, a Ty bierzesz pierścionek i 2 cegły. Jeśli Ci pasuje, złożę taką propozycję.} \\
            }
        }
    }
    & & \\ \\

    \theutterance \stepcounter{utterance}  
    & & & \multicolumn{2}{p{0.3\linewidth}}{
        \cellcolor[rgb]{0.9,0.9,0.9}{
            \makecell[{{p{\linewidth}}}]{
                \texttt{\tiny{[GM$|$GM]}}
                \texttt{continue} \\
            }
        }
    }
    & & \\ \\

    \theutterance \stepcounter{utterance}  
    & & & \multicolumn{2}{p{0.3\linewidth}}{
        \cellcolor[rgb]{0.9,0.9,0.9}{
            \makecell[{{p{\linewidth}}}]{
                \texttt{\tiny{[GM$|$GM]}}
                \texttt{Brzmi dobrze. Żeby domknąć to bez ryzyka konfliktu propozycji, proponuję dokładnie tak: ja biorę zegarek i obcęgi, a Ty bierzesz pierścionek i 2 cegły. Jeśli Ci pasuje, złożę taką propozycję.} \\
            }
        }
    }
    & & \\ \\

    \theutterance \stepcounter{utterance}  
    & & \multicolumn{4}{p{0.6\linewidth}}{
        \cellcolor[rgb]{0.9,0.9,0.9}{
            \makecell[{{p{\linewidth}}}]{
                \texttt{\tiny{[GM$\rangle$P2]}}
                \texttt{Brzmi dobrze. Żeby domknąć to bez ryzyka konfliktu propozycji, proponuję dokładnie tak: ja biorę zegarek i obcęgi, a Ty bierzesz pierścionek i 2 cegły. Jeśli Ci pasuje, złożę taką propozycję.} \\
            }
        }
    }
    & \\ \\

    \theutterance \stepcounter{utterance}  
    & & & \multicolumn{4}{p{0.6\linewidth}}{
        \cellcolor[rgb]{0.9,0.9,0.9}{
            \makecell[{{p{\linewidth}}}]{
                \texttt{\tiny{[GM$\langle$P2]}}
                \texttt{{[}Propozycja: 1 pierścionek, 2 cegły{]}} \\
            }
        }
    }
     \\ \\

    \theutterance \stepcounter{utterance}  
    & & & \multicolumn{2}{p{0.3\linewidth}}{
        \cellcolor[rgb]{0.9,0.9,0.9}{
            \makecell[{{p{\linewidth}}}]{
                \texttt{\tiny{[GM$|$GM]}}
                \texttt{proposal} \\
            }
        }
    }
    & & \\ \\

    \theutterance \stepcounter{utterance}  
    & & & \multicolumn{2}{p{0.3\linewidth}}{
        \cellcolor[rgb]{0.9,0.9,0.9}{
            \makecell[{{p{\linewidth}}}]{
                \texttt{\tiny{[GM$|$GM]}}
                [0, 2, 1, 0]
            }
        }
    }
    & & \\ \\

    \theutterance \stepcounter{utterance}  
    & & \multicolumn{4}{p{0.6\linewidth}}{
        \cellcolor[rgb]{0.9,0.9,0.9}{
            \makecell[{{p{\linewidth}}}]{
                \texttt{\tiny{[P1$\langle$GM]}}
                \texttt{Drugi gracz złożył tajną propozycję. Teraz Twoja kolej na złożenie własnej propozycji.} \\
            }
        }
    }
    & \\ \\

    \theutterance \stepcounter{utterance}  
    & \multicolumn{4}{p{0.6\linewidth}}{
        \cellcolor[rgb]{0.9,0.9,0.9}{
            \makecell[{{p{\linewidth}}}]{
                \texttt{\tiny{[P1$\rangle$GM]}}
                \texttt{{[}Propozycja: 1 zegarek, 1 obcęgi{]}} \\
            }
        }
    }
    & & \\ \\

    \theutterance \stepcounter{utterance}  
    & & & \multicolumn{2}{p{0.3\linewidth}}{
        \cellcolor[rgb]{0.9,0.9,0.9}{
            \makecell[{{p{\linewidth}}}]{
                \texttt{\tiny{[GM$|$GM]}}
                \texttt{proposal} \\
            }
        }
    }
    & & \\ \\

    \theutterance \stepcounter{utterance}  
    & & & \multicolumn{2}{p{0.3\linewidth}}{
        \cellcolor[rgb]{0.9,0.9,0.9}{
            \makecell[{{p{\linewidth}}}]{
                \texttt{\tiny{[GM$|$GM]}}
                [1, 0, 0, 1]
            }
        }
    }
    & & \\ \\

    \theutterance \stepcounter{utterance}  
    & & & \multicolumn{2}{p{0.3\linewidth}}{
        \cellcolor[rgb]{0.9,0.9,0.9}{
            \makecell[{{p{\linewidth}}}]{
                \texttt{\tiny{[GM$|$GM]}}
                [[1, 0, 0, 1], [0, 2, 1, 0]]
            }
        }
    }
    & & \\ \\

\end{supertabular}
}

\end{document}
